\documentclass[aps,prb,onecolumn,superscriptaddress]{revtex4-2}
\usepackage{amsmath,amssymb,graphicx,bm}
\usepackage{underscore}
\usepackage{color}
\usepackage[colorlinks=true, linkcolor=blue, citecolor=blue, urlcolor=blue]{hyperref}
%\usepackage{bibunits}  % For separate SI references

\renewcommand{\thefigure}{S\arabic{figure}}
\renewcommand{\thetable}{S\arabic{table}}
\renewcommand{\theequation}{S\arabic{equation}}

\makeatletter
\renewcommand{\@biblabel}[1]{[S#1]}
\makeatother

% \usepackage[labelformat=empty, justification=raggedright]{caption}

\begin{document}
% \begin{bibunit}[apsrev4-2]  % Start a separate reference unit for SI

\title{Supplemental Material for: \\ \textit{Synchronization and Spatiotemporal Phases in Driven Precessing Domain Walls}}
\maketitle
\section{Movies}

\begin{itemize}
\item \textbf{Movie S1} (\texttt{movie\_StabilityMapVsGilbertDamping.mp4}) illustrates how the spectral stability map evolves as a function of the Gilbert damping parameter $\alpha$. In each frame, color indicates the magnitude of the largest Floquet multiplier for each mode at a given field value.
\item \textbf{Movie S2} (\texttt{movie\_BlochLineDynamics.mp4}) shows the space-time trajectories of vertical Bloch lines within a segment of a domain wall embedded in a larger DW structure. The colorbar represents the instantaneous angle $\phi(x,t)$, with Bloch lines appearing as kinks that generate the spatiotemporal pattern.
\item \textbf{Movie S3} (\texttt{movie\_DWBlochLinesCorrugationLLG.mp4}) 
Emergence of corrugations and Bloch lines in the transient domain-wall dynamics, starting from a slightly perturbed flat initial condition, obtained from LLG numerical simulations.
\end{itemize}
\end{document}